\chapter{Event Selection}
\section{Online}
\subsection{Trigger Eras}
%\subsection{Electron Trigger}
%\subsubsection{Electron Algorithms}
%\subsubsection{Trigger Performance}
\subsection{Level-1 Trigger}
\subsubsection{Description of Algorithm}
Description of algorithm.  
\subsubsection{Performance}
Distributions, resolutions, efficiencies of Level-1 electron/photon objects.  

Figures: Trigger Electron/Photon Pt, eta, phi spectra for L1, Fig. \ref{fig:L1TriggerObjectSpectra}

 \begin{figure}[htb]
  \begin{center}
    \includegraphics[width=360pt]{CMS-BW.pdf}
  \end{center}
  \caption[Trigger Electron/Photon Pt, eta, phi spectra for L1]{Trigger Electron/Photon Pt, eta, phi spectra for L1.}
  \label{fig:L1TriggerObjectSpectra}
 \end{figure}


Figures: \DR for L1 trigger object matching, Fig. \ref{fig:L1TriggerObjectDeltaR}

 \begin{figure}[htb]
  \begin{center}
    \includegraphics[width=360pt]{CMS-BW.pdf}
  \end{center}
  \caption[\DR for L1 trigger object matching to offline]{\DR for L1 trigger object matching to offline-reconstructed objects.}
  \label{fig:L1TriggerObjectDeltaR}
 \end{figure}

Figures: L1 trigger Pt resolution plots using offline reference, Fig. \ref{fig:L1TriggerObjectResolutions}

 \begin{figure}[htb]
  \begin{center}
    \includegraphics[width=360pt]{CMS-BW.pdf}
  \end{center}
  \caption[L1 trigger Pt resolution plots using offline reference]{L1 trigger Pt resolution plots using offline reference.}
  \label{fig:L1TriggerObjectResolutions}
 \end{figure}

Figures: L1 trigger electron/photon efficiency turn on using offline reconstructed electrons and MC, Fig. \ref{fig:L1TriggerObjectEfficiencies}

 \begin{figure}[htb]
  \begin{center}
    \includegraphics[width=360pt]{CMS-BW.pdf}
  \end{center}
  \caption[L1 trigger electron/photon efficiency turn on using offline reconstructed electrons and MC]{L1 trigger electron/photon efficiency turn on using offline reconstructed electrons and MC.}
  \label{fig:L1TriggerObjectEfficiencies}
 \end{figure}


\subsection{High-Level Trigger}
\subsubsection{Description of Algorithms}
Description of algorithms.
\subsubsection{Performance}
Distributions, resolutions, efficiencies of HLT electron/photon objects.

Figures: Trigger Electron/Photon Pt, eta, phi spectra for HLT, Fig. \ref{fig:HLTriggerObjectSpectra}

 \begin{figure}[htb]
  \begin{center}
    \includegraphics[width=360pt]{CMS-BW.pdf}
  \end{center}
  \caption[Trigger Electron/Photon Pt, eta, phi spectra for HLT]{Trigger Electron/Photon Pt, eta, phi spectra for HLT.}
  \label{fig:HLTriggerObjectSpectra}
 \end{figure}


Figures: \DR for HLT trigger object matching, Fig. \ref{fig:HLTriggerObjectDeltaR}

 \begin{figure}[htb]
  \begin{center}
    \includegraphics[width=360pt]{CMS-BW.pdf}
  \end{center}
  \caption[\DR for HLT trigger object matching to offline]{\DR for HLT trigger object matching to offline-reconstructed objects.}
  \label{fig:HLTriggerObjectDeltaR}
 \end{figure}

Figures: HLT trigger Pt resolution plots using offline reference, Fig. \ref{fig:HLTriggerObjectResolutions}

 \begin{figure}[htb]
  \begin{center}
    \includegraphics[width=360pt]{CMS-BW.pdf}
  \end{center}
  \caption[HLT trigger Pt resolution plots using offline reference]{HLT trigger Pt resolution plots using offline reference.}
  \label{fig:HLTriggerObjectResolutions}
 \end{figure}

Figures: HLT trigger electron/photon efficiency turn on using offline reconstructed electrons and MC, Fig. \ref{fig:HLTriggerObjectEfficiencies}

 \begin{figure}[htb]
  \begin{center}
    \includegraphics[width=360pt]{CMS-BW.pdf}
  \end{center}
  \caption[HLT trigger electron/photon efficiency turn on using offline reconstructed electrons and MC]{HLT trigger electron/photon efficiency turn on using offline reconstructed electrons and MC.}
  \label{fig:HLTriggerObjectEfficiencies}
 \end{figure}




\section{Offline}
\subsection{Electron Selection}
Starts out with full set of objects reconstructed as electrons.  
NOT necessarily real electrons, but energy deposits consistent with those from electrons at a low-level.  
(talk about def of GSF electron etc so know what kinds of cuts have already been applied)
Further conditions are applied to select the electron objects most likely to be real electrons from a Z.  
First electrons are matched to trigger objects passing an egamma trigger with a Level-1 threshold of 5 GeV.  
A 20-GeV cut on the supercluster transverse energy is applied.  
Next comes a set of cuts designed to reject electrons from photon conversions.  
Then come cuts to distinguish real electrons from other energy deposits which may mimic electrons, 
and finally are cuts to make sure the electron is isolated, i.e. not part of a jet (which would not be coming from a Z).  

\subsubsection{Matching Offline-Reconstructed Objects to Trigger Objects}
Figures: \DR for L1 and HLT trigger object matching, Fig. \ref{fig:TriggerObjectSelectionDeltaR}

 \begin{figure}[htb]
  \begin{center}
    \includegraphics[width=360pt]{CMS-BW.pdf}
  \end{center}
  \caption[\DR for L1 and HLT trigger object matching to offline for electron selection]{\DR for L1 and HLT trigger object matching to offline-reconstructed objects.}
  \label{fig:TriggerObjectSelectionDeltaR}
 \end{figure}


\subsubsection{Conversion Rejection}
The conversion rejection step of the electron selection aims to eliminate photons that have converted into electron pairs.  
Several variables have been chosen to distinguish between conversions and prompt electrons, focusing on the track properties.  
Since prompt electrons are expected to hit the first layer of the tracker, 
any electron-associated track that is missing initial hits can be rejected as a conversion that occurred within the tracker.  
The looser requirements of the 95 percent efficient working point allow an electron to have one missing hit;
tighter working points require zero missing hits.  
Fig. \ref{fig:MissHitsConvRejVars} shows the number of missing hits for each electron.  
In addition, the conversion rejection algorithm searches for candidate electron-pair partner tracks, 
i.e. track pairs that appear to come from the same point.  
Every electron track is compared to all opposite-charge tracks within a dR cone of 0.3, and two values are calculated: 
a measure of the difference in the angles in the r-z plane between the two tracks (dCotTheta, Fig. \ref{fig:DCotConvRejVars}), 
and the distance of closest approach between the two tracks in the x-y plane (dist, Fig. \ref{fig:DistConvRejVars}).  
Electron pairs from conversions are expected to have low values for both quantities, 
since both electron tracks begin at the point of conversion and initially have the same direction as the converting photon.  
The WP95 requirements apply a threshold at 0.02 for both quantities.  
If an electron track corresponds to a track pair having values below that threshold for both variables, 
it is assumed to be a conversion and is rejected.  

%Figures: Conversion rejection variables before respective cuts, Fig. \ref{fig:ConvRejVars}

 \begin{figure}[htb]
  \begin{center}
%    \includegraphics[width=180pt, angle=90]{missHitsSeq-14sep10.pdf}
    \includegraphics[width=180pt, angle=90]{missHitsSeq.pdf}
  \end{center}
  \caption[Missing hits conversion rejection variables after previous cuts]{Missing hits conversion rejection variables after previous cuts.}
  \label{fig:MissHitsConvRejVars}
 \end{figure}

 \begin{figure}[htb]
  \begin{center}
%    \includegraphics[width=180pt, angle=90]{dCotThetaSeq-14sep10.pdf}
    \includegraphics[width=180pt, angle=90]{dCotThetaSeq.pdf}
  \end{center}
  \caption[DCotTheta conversion rejection variables after previous cuts]{DCotTheta conversion rejection variables after previous cuts.}
  \label{fig:DCotConvRejVars}
 \end{figure}

 \begin{figure}[htb]
  \begin{center}
%    \includegraphics[width=180pt, angle=90]{distSeq-14sep10.pdf}
    \includegraphics[width=180pt, angle=90]{distSeq.pdf}
  \end{center}
%  \caption[Conversion rejection variables before respective cuts]{Conversion rejection variables before respective cuts.}
  \caption[Dist conversion rejection variables after previous cuts]{Dist conversion rejection variables after previous cuts.}
  \label{fig:DistConvRejVars}
 \end{figure}


\subsubsection{Electron Identification}
Several variables to differentiate genuine electrons from other particles that may mimic electron behavior.  

Figures: Electron ID variables before respective cuts, Fig. \ref{fig:ElecIdVars}

 \begin{figure}[htb]
  \begin{center}
    \includegraphics[width=360pt]{CMS-BW.pdf}
  \end{center}
  \caption[Electron ID variables before respective cuts]{Electron ID variables before respective cuts.}
  \label{fig:ElecIdVars}
 \end{figure}



Figures: Reconstructed electron/photon Pt, eta, phi spectra after ID, Fig. \ref{fig:RecoSpectraAfterEid}

 \begin{figure}[htb]
  \begin{center}
    \includegraphics[width=360pt]{CMS-BW.pdf}
  \end{center}
  \caption[Reconstructed electron/photon Pt, eta, phi spectra after ID]{Reconstructed electron/photon Pt, eta, phi spectra after ID.}
  \label{fig:RecoSpectraAfterEid}
 \end{figure}


\subsubsection{Electron Isolation}
Conditions to require that electrons found are isolation, i.e. not contained in a jet with other particles.  
We want (``prompt''?) electrons that come directly from the Z, not from some other interaction with lots of other end products packaged along with the electron.  

Figures: Electron isolation variables before respective cuts, Fig. \ref{fig:ElecIsoVars}

 \begin{figure}[htb]
  \begin{center}
    \includegraphics[width=360pt]{CMS-BW.pdf}
  \end{center}
  \caption[Electron isolation variables before respective cuts]{Electron isolation variables before respective cuts.}
  \label{fig:ElecIsoVars}
 \end{figure}



Figures: Reconstructed electron/photon Pt, eta, phi spectra after ID + isolation, Fig. \ref{fig:RecoSpectraAfterEidEiso}

 \begin{figure}[htb]
  \begin{center}
    \includegraphics[width=360pt]{CMS-BW.pdf}
  \end{center}
  \caption[Reconstructed electron/photon Pt, eta, phi spectra after ID + isolation]{Reconstructed electron/photon Pt, eta, phi spectra after ID + isolation.}
  \label{fig:RecoSpectraAfterEidEiso}
 \end{figure}


\subsubsection{Electron Selection Performance} %wesley comment
Figures: Reco PT resolution plots using MC reference for Z MC, Fig. \ref{fig:RecoPtResolution}

 \begin{figure}[htb]
  \begin{center}
    \includegraphics[width=360pt]{CMS-BW.pdf}
  \end{center}
  \caption[Reco PT resolution plots using MC reference for Z MC]{Reco PT resolution plots using MC reference for Z MC.}
  \label{fig:RecoPtResolution}
 \end{figure}


\subsection{\Zee Event Selection}
\subsubsection{Acceptance for \Zee Events}
\subsubsection{Efficiency of Selection for \Zee Events}
Tag and Probe:  The tag and probe method is used to determine the efficiency of a given set of electron selection steps.  
In this method, a set of data events is identified which appear to be good Z->ee candidate events.  
One of the electrons in each of these events must pass 
all selection cuts and essentially be identified as a golden electron; 
%some stringent
this is the tag.  
%The other electron must 
%pass some less stringent criteria
%be identified at least as a supercluster; 
The criteria applied to the other electron are looser: it is not required to pass the selection cuts being tested.  
This is the probe electron.  
The event's status as a Z candidate is determined by enforcing a mass window cut on the invariant mass of the tag-probe system 
(this is sufficient(enough?) to select events in which both legs are real electrons -- REFERENCE??).  
The selection efficiency is then given by the fraction of probe electrons that pass the requirement being tested.  

Figures: Efficiency table and supporting plots

\subsubsection{Distribution of Z Kinematic Variables}
